\documentclass[11pt]{article}
\usepackage[OT1]{fontenc}
\renewcommand{\rmdefault}{cmr}
\renewcommand{\sfdefault}{cmss}
\renewcommand{\ttdefault}{cmtt}
\usepackage{amsmath,amssymb,amsthm}
\usepackage[hidelinks]{hyperref}
\setlength{\textwidth}{6.1in}
\setlength{\oddsidemargin}{0.2in}
\setlength{\evensidemargin}{0.2in}
\setlength{\textheight}{9in}
\setlength{\topmargin}{-0.35in}
\setlength{\headheight}{13.6pt}
\setlength{\emergencystretch}{1.5em}
\newtheorem{theorem}{Theorem}[section]
\newtheorem{lemma}[theorem]{Lemma}
\newtheorem{proposition}[theorem]{Proposition}
\newtheorem{corollary}[theorem]{Corollary}
\theoremstyle{definition}
\newtheorem{definition}[theorem]{Definition}
\newtheorem{example}[theorem]{Example}
\newcommand{\Comp}{\operatorname{Comp}}
\newcommand{\Forb}{\operatorname{Forb}_{\rm ind}}
\newcommand{\E}{\mathsf E}
\newcommand{\C}{\mathcal C}
\newcommand{\B}{\mathcal B}
\newcommand{\OO}{\mathcal O}
\newcommand{\NN}{\operatorname{N}}
\newcommand{\dSigma}{d\text{-}\Sigma}

\title{Scott complexity of countable abelian $p$-groups:\\classification in the reduced case and the realization spectrum}
\author{}
\date{21 October 2026\\Working manuscript}
\begin{document}
\maketitle
\begin{abstract}
We determine the exact Scott complexity of every countable reduced
abelian $p$-group. Write its length as $\omega\gamma+n$, $n<\omega$.
At a limit length the complexity is $\Pi_{2\gamma+1}$; otherwise the
bounded terminal subgroup at $\omega\gamma$ determines one of four
exact complexities. The distinction depends on its infinite cyclic
multiplicities and on shorter finite summands, rather than on how many
infinite Ulm invariants occur anywhere in the group. We construct reduced
examples at every positive countable successor difference-of-$\Sigma$
level. We also exclude exact pure $\Sigma$ complexity for every countable
abelian $p$-group, including those with divisible summands. This answers
Questions~5.1 and~5.2 of Alvir--Csima--MacLean in their stated versions.
Their reduced-group back-and-forth theorem, Ulm classification, Zippin
existence, and general Scott-complexity restrictions are credited inputs.
The classification for individual nonreduced groups remains open here.
\end{abstract}

\section{The problem, its source, and a first example}
A Scott sentence describes one countable structure up to isomorphism
among countable models. Its complexity measures how hard that description
must be in countable infinitary logic. Algebraic classification alone does
not settle the optimal logical complexity: it supplies descriptions, but
one must still show that simpler descriptions cannot work.

Alvir, Csima and MacLean \cite{ACM} give upper bounds and extensive
back-and-forth analysis for reduced abelian $p$-groups. Their
Question~5.1 asks about optimality when at least two Ulm invariants are
infinite. Their Question~5.2 asks which pure $\Sigma$ and difference
complexities occur. Gonzalez, Rossegger and Turetsky \cite{GRT} later
ask for explicit abelian $p$-groups displaying a gap between two Scott
rank conventions and for the non-$\Pi$ spectrum. Their abstract
existence argument is prior work, rather than a new conclusion here.

The simplest useful example is
\[
 C_p\oplus C_{p^2}^{(\omega)}.
\]
It has exact $d\text{-}\Sigma_2$ complexity. The finite $C_p$ summand
can be described existentially, but it can disappear in a suitable
increasing union of copies of the group, obstructing a $\Pi_2$ Scott
sentence. If the $C_p$ summand is absent, the complexity drops to
$\Pi_2$. If both multiplicities are infinite, it rises to $\Pi_3$.
The same three-way distinction reappears at any terminal
$\omega$-block of the transfinite $p$-filtration, with a separate
finite-terminal case. The classification below makes that assertion exact.

The proof has three main steps. First we give the bounded descriptions.
Next we construct quotients which change a terminal invariant while
preserving specified back-and-forth comparisons. The same quotient must
serve both comparisons in the highest-complexity case; separate
competitors would not exclude a difference sentence. Finally we build
explicit transfinite examples and handle the divisible part separately.

\section{Conventions and the classification}
Fix a prime $p$. Groups are countable abelian $p$-groups in the original
language $\{+,-,0\}$. No filtration predicates are added to that language.
Write $C_{p^j}$ for the cyclic group of order $p^j$ and
$C_{p^\infty}$ for the Pr\"ufer group. The standard decomposition is
$G=R\oplus D$, with $R$ reduced and $D$ divisible. Define
\[
 G_0=G,\qquad G_{\xi+1}=pG_\xi,\qquad
 G_\lambda=\bigcap_{\xi<\lambda}G_\xi.
\]
\[
 P_\xi(G)=G_\xi[p],\qquad
 u_G(\xi)=\dim_{\mathbb F_p}P_\xi(G)/P_{\xi+1}(G).
\]
The reduced length $\lambda_R$ is the least index with $R_{\lambda_R}=0$.
For $x\ne0$, its height is $\xi$ when $x\in G_\xi\setminus G_{\xi+1}$;
elements in the divisible part and zero have height $\infty$.
All ordinal products are ordinary ordinal products:
$2\gamma=2\cdot\gamma$ and $\omega\gamma=\omega\cdot\gamma$.

We use ordinary, not necessarily computable, infinitary
$\Sigma_\alpha$ and $\Pi_\alpha$ formulas. At level zero they are
quantifier-free; a $\Sigma_\alpha$ formula is a countable disjunction
of existential formulas with $\Pi_\beta$ matrices, $\beta<\alpha$,
and $\Pi_\alpha$ is the dual class. A $d\text{-}\Sigma_\alpha$
sentence is a conjunction of a $\Sigma_\alpha$ and a $\Pi_\alpha$
sentence. \emph{Exact} complexity refers to the usual Scott-complexity
hierarchy; it does not deny that a structure has padded descriptions at
higher levels. We use the general restrictions and Miller collapse
from \cite[Theorems~1.6, 2.5 and 2.6]{AGHT}.

For finite tuples, $(A,\bar a)\leq_0(B,\bar b)$ means atomic agreement.
For $\alpha>0$, every $\beta<\alpha$ and finite right extension $\bar d$
must have a left response $\bar c$ with
$(B,\bar b\bar d)\leq_\beta(A,\bar a\bar c)$.
Thus $\Pi_\alpha$ truth transfers from the left to the right,
and $\Sigma_\alpha$ truth transfers from the right to the left.
This direction will be used throughout.

\begin{theorem}[Exact reduced-group classification]
\label{thm:reduced-classification}
Let $G\ne0$ be reduced, and write its length uniquely as
$\lambda=\omega\gamma+n$, where $n<\omega$.
If $n=0$, then $\gamma>0$ and $G$ has exact complexity
$\Pi_{2\gamma+1}$.
If $n>0$, put $H=G_{\omega\gamma}$ and write
\[
 H\cong\bigoplus_{j=1}^{n}C_{p^j}^{(m_j)},\qquad
 m_j\in\mathbb N\cup\{\omega\},\quad m_n>0.
\]
The exact complexity is
\begin{center}
\normalfont
\begin{tabular}{@{}p{0.62\linewidth}l@{}}
\hline
Condition on $H$ & Exact complexity\\
\hline
$H$ is finite & $d\text{-}\Sigma_{2\gamma+1}$\\[4pt]
Exactly one $m_k=\omega$, and $m_j=0$ for every $j<k$
 & $\Pi_{2\gamma+2}$\\[4pt]
Exactly one $m_k=\omega$, and $m_j>0$ for some $j<k$
 & $d\text{-}\Sigma_{2\gamma+2}$\\[4pt]
At least two $m_j=\omega$ & $\Pi_{2\gamma+3}$\\
\hline
\end{tabular}
\end{center}
The trivial group has exact complexity $\Pi_1$.
\end{theorem}


The terminal multiplicities are those of the bounded group
$G_{\omega\gamma}$, not the complete list of Ulm invariants of $G$.
For example, finite summands above its unique infinite multiplicity are
allowed in both middle rows. Their presence does not change which row
applies.

\section{The bounded base case}
\begin{theorem}\label{thm:bounded-pgroup-classification}
Let $G$ be a countably infinite Abelian $p$-group of bounded exponent,
and write
\[
 G\cong\bigoplus_{j=1}^n C_{p^j}^{(m_j)},
 \qquad m_j\in\mathbb N\cup\{\omega\}.
\]
Its Scott complexity is determined as follows.
\begin{center}
\begin{tabular}{p{0.67\linewidth}|l}
Condition & Exact complexity\\ \hline
At least two $m_j$ equal $\omega$ & $\Pi_3$\\
Exactly one $m_k=\omega$, and some $m_j>0$ with $j<k$
 & $d\text{-}\Sigma_2$\\
Exactly one $m_k=\omega$, and all $m_j=0$ with $j<k$
 & $\Pi_2$
\end{tabular}
\end{center}
Thus no bounded-exponent countable Abelian $p$-group has an exact
pure $\Sigma$ Scott complexity. The case of at least two infinite
multiplicities is \cite[Theorem~4.2]{ACM}; the proof below supplies
the remaining cases, including arbitrary finite summands above the
unique infinite multiplicity.
\end{theorem}

\begin{proof}
The decomposition is the standard bounded-exponent case of the Ulm
classification. Suppose that the only infinite multiplicity is at $k$,
and write
\[
 G=F_-\oplus C_{p^k}^{(\omega)}\oplus F_+,
\]
where $F_-$ has cyclic summands of orders below $p^k$, and $F_+$ has
cyclic summands of orders above $p^k$. Both $F_-$ and $F_+$ are finite.
Let $K=p^kF_+$, a fixed finite group.

The condition $p^kG\cong K$ has a $d\text{-}\Sigma_1$ description.
If $K$ has $M$ elements, an existential tuple of elements
$p^ky_1,\ldots,p^ky_M$ can carry its complete finite group table,
including all required inequalities. This is $\Sigma_1$.
The bound $|p^kG|\le M$ is $\Pi_1$, expressed by
\[
 \forall y_0\cdots y_M\,
 \bigvee_{0\le i<j\le M}p^ky_i=p^ky_j.
\]
Together these conditions fix all cyclic multiplicities above $k$.
The condition that $p^{k-1}G$ is infinite is $\Pi_2$:
\[
 \bigwedge_{m\ge1}\exists y_1\cdots y_m
 \bigwedge_{i<j}p^{k-1}y_i\ne p^{k-1}y_j.
\]
Together with finiteness of $p^kG$, it forces the multiplicity at
$k$ to be $\omega$.

For each $j<k$, the specified finite multiplicity $m_j$ is a finite
Ulm invariant. The usual lower dimension condition is $\Sigma_2$,
and the corresponding upper condition is $\Pi_2$: choose elements in
$(p^{j-1}G)[p]$ and require their nonzero linear combinations to lie
outside $p^jG$. A finite conjunction of these exact dimension
conditions is $d\text{-}\Sigma_2$.
Add the Abelian group axioms and the bound $p^nG=0$.
These conditions determine the entire cyclic decomposition, and hence
give a $d\text{-}\Sigma_2$ Scott sentence.

If $F_-=0$, replace all the lower dimension conditions by
\[
 \forall x\left(px=0\ \longrightarrow\
                 \exists y\,p^{k-1}y=x\right).
\]
This $\Pi_2$ sentence says that every socle element has height at least
$k-1$, and so there are no cyclic summands shorter than $p^k$.
The remaining conditions above are all $\Pi_2$ or of lower complexity.
Thus $G$ has a $\Pi_2$ Scott sentence in this case. It is exact:
an infinite group cannot have a universal Scott sentence because such
a sentence passes to the trivial subgroup; a $d\text{-}\Sigma_1$
sentence would pass to the finite subgroup generated by its finite
existential witnesses.

If $F_-\ne0$, no $\Sigma_2$ Scott sentence is possible. More
generally, no infinite locally finite group has a $\Sigma_2$ Scott
sentence. One true existential disjunct has finite witnesses and a
$\Pi_1$ matrix. The witnesses generate a finite subgroup, and the
matrix passes to that subgroup. The purported Scott sentence would
therefore have a nonisomorphic finite countable model.

To exclude $\Pi_2$ in the case $F_-\ne0$, choose cyclic generators
$u_1,\ldots,u_t$ for $F_-$ with orders $p^{\ell_1},\ldots,p^{\ell_t}$,
where $1\le\ell_j<k$, and generators $e_0,e_1,\ldots$ for the
$C_{p^k}$ summands. Define an injective endomorphism by
\[
 f(u_j)=p^{k-\ell_j}e_{j-1},\qquad
 f(e_i)=e_{i+t},\qquad f\upharpoonright F_+=\operatorname{id}.
\]
Orders are preserved and the specified summand images have disjoint
supports. In its direct limit, the finite $F_+$ factor stays fixed.
The remaining factor has exponent $p^k$, an infinite
$p^{k-1}$-multiple subgroup, and every socle element has a
$p^{k-1}$-root at a later stage: one application of $f$ moves the socle
of each short summand into $p^{k-1}\langle e_{j-1}\rangle$.
The cyclic decomposition therefore makes the limit
\[
 D\cong C_{p^k}^{(\omega)}\oplus F_+,
\]
which is nonisomorphic to $G$ because its lower multiplicities are all
zero. A $\Pi_2$ sentence is preserved in an increasing union of
substructures: every finite universal tuple occurs at a stage, whose
finite existential witnesses and quantifier-free diagram remain in the
union. Thus no $\Pi_2$ Scott sentence for $G$ exists. The
$d\text{-}\Sigma_2$ upper bound is consequently exact.

The only remaining infinite case has at least two infinite
multiplicities, whose exact $\Pi_3$ complexity is the cited theorem of
Alvir--Csima--MacLean. Finite nontrivial groups have exact
$d\text{-}\Sigma_1$, and the trivial group has a $\Pi_1$ Scott sentence,
so finite groups add no pure $\Sigma$ exception.
\end{proof}


\section{Changing tails while controlling comparisons}
We use Ulm classification and the comparison theorem
\cite[Theorem~4.4]{ACM} for reduced groups. Its odd and even cases
include alternatives expressed by relative $\Sigma_1$ and $\Sigma_2$
truth in the terminal subgroup. We check those alternatives explicitly;
they are not discarded when a terminal socle is finite.

A second input is Zippin's existence theorem, in the form stated by
Calderoni and Thomas \cite[Theorem~4.2]{CT}: a countable sequence of
nonzero countable direct sums of cyclic $p$-groups is the Ulm-factor
sequence of a reduced group exactly when every nonfinal factor has
unbounded exponent. Here the factors are
$U_\xi(G)=G_{\omega\xi}/G_{\omega(\xi+1)}$.
We use that published theorem rather than a proposed infinite reverse
iteration of one-block constructions.

\begin{lemma}[Transfinite prescribed-tail existence]
\label{lem:prescribed-tail}
Let $\gamma>0$ be countable, put $\delta=\omega\gamma$, and let $G$
be a countable reduced Abelian $p$-group of length at least $\delta$.
For every nonzero countable reduced $p$-group $T$ there is a countable
reduced $p$-group $H(T)$ such that
\[
 u_{H(T)}(\xi)=u_G(\xi)\quad(\xi<\delta),
 \qquad H(T)_\delta\cong T.
\]
Its length is $\delta+\lambda_T$.
\end{lemma}

\begin{proof}
For each $\xi<\gamma$, $G_{\omega\xi}$ has length at least $\omega$.
Its nonzero finite Ulm invariants have unbounded support: otherwise a
finite-stage subgroup would be divisible, contradicting reducedness.
Consequently every $U_\xi(G)$ with $\xi<\gamma$ is nonzero and has
unbounded exponent, including the last one when $\gamma$ is a successor.

Let $\tau(T)$ be the ordinal length of the nonzero Ulm-factor sequence
of $T$. Concatenate the first $\gamma$ Ulm factors of $G$ with the entire Ulm
factor sequence of $T$. This is a countable sequence indexed by
$\gamma+\tau(T)$. Every nonfinal factor is unbounded: this has just
been checked for the initial factors, and it holds for the nonfinal
factors of the reduced group $T$. Zippin's theorem supplies a reduced
group $H$ with this sequence. The Ulm factors of $H_\delta$ are exactly
those of $T$, so Ulm's classification gives $H_\delta\cong T$.
The initial multiplicities give the displayed lower Ulm identities.
Finally $H_{\delta+\eta}=(H_\delta)_\eta$ and no earlier initial
factor vanishes, giving length $\delta+\lambda_T$.
\end{proof}

\begin{lemma}[Quotient filtration below a limit]
\label{lem:quotient-filtration}
Let $\delta=\omega\gamma$ and $T\le G_\delta$. For the quotient
map $q:G\to Q=G/T$, one has
\[
 Q_\alpha=q(G_\alpha)\quad(\alpha\le\delta).
\]
For every $\beta<\delta$, $q$ induces an isomorphism of the
$\beta$-th Ulm quotients, so $u_Q(\beta)=u_G(\beta)$.
\end{lemma}

\begin{proof}
The successor step follows from $p q(G_\alpha)=q(pG_\alpha)$.
At a limit $\eta\le\delta$, the subgroup $T$ lies in every
$G_\alpha$, $\alpha<\eta$, so
$q^{-1}(q(G_\alpha))=G_\alpha$. Intersecting preimages proves
the limit step, including $\eta=\delta$.

For the Ulm quotient, $\beta+2<\delta$ whenever $\beta<\delta$;
the case $\delta=0$ has no such $\beta$. Let $y\in P_\beta(Q)$
and choose $x\in G_\beta$ with $q(x)=y$. Then
$px\in T\subseteq G_{\beta+2}$, so there is $z\in G_{\beta+1}$
with $pz=px$. Thus $x-z\in P_\beta(G)$ and
$q(x-z)$ differs from $y$ by the element $q(z)\in P_{\beta+1}(Q)$.
This gives surjectivity on Ulm quotients. If $x\in P_\beta(G)$ and
$q(x)\in P_{\beta+1}(Q)$, then
$x\in G_{\beta+1}+T=G_{\beta+1}$, proving injectivity.
\end{proof}


\begin{lemma}[Finite-product preservation]\label{lem:products}
If $(A_i,\bar a_i)\leq_\alpha(B_i,\bar b_i)$ for $i=1,2$, then
\[
 (A_1\oplus A_2,(\bar a_1,\bar a_2))
 \leq_\alpha
 (B_1\oplus B_2,(\bar b_1,\bar b_2)).
\]
The paired notation denotes componentwise pairing of corresponding
entries into tuples of product elements.
\end{lemma}
\begin{proof}
Induct on $\alpha$. Equations between group terms in the product hold
exactly when they hold in both coordinates, so atomic agreement is
preserved. Given a right-hand finite extension and $\beta<\alpha$,
split it into its coordinate tuples. Apply the two given comparisons
to obtain coordinate responses. The induction hypothesis at $\beta$
combines their reversed comparisons into the response in the product.
This proves the assertion at successors and limits alike.
\end{proof}

\subsection{The obstruction at odd-form levels}
The explicit groups $B_\omega,W_\omega$ used next are defined and proved
to have the indicated filtrations in Section~\ref{sec:walker}.
\begin{theorem}
\label{thm:odd-obstruction}
Let $\gamma>0$ be countable. If $G=R\oplus D$, with $R$ reduced and
$D$ divisible, and $\lambda_R\ge\omega\gamma$, then $G$ has no
$\Sigma_{2\gamma+1}$ Scott sentence. No countable Abelian
$p$-group has exact pure $\Sigma_{2\gamma+1}$ complexity.
\end{theorem}

\begin{proof}
Put $\delta=\omega\gamma$. Take the two reduced tails
\[
 T_0=\bigoplus_{n<\omega}C_{p^{n+1}}^{(\omega)},
 \qquad T_1=B_\omega\oplus W_\omega,
\]
where $B_\omega,W_\omega$ are the presentations in Section~\ref{sec:walker}.
Their lengths are respectively $\omega$ and $\omega+1$, and their
socles at every finite height are infinite. Lemma~\ref{lem:prescribed-tail}
gives $H_i=H(T_i)$ with the same Ulm invariants as $R$ below $\delta$
and with distinct lengths $\delta+\omega$, $\delta+\omega+1$.

Apply ACM, Theorem~4.4(1) \cite{ACM}, with left group $H_i$, right group $R$,
index $\gamma$, and empty tuples. Condition (a) is the equality of
lower invariants. Condition (b) holds because
\[
 \lambda_{H_i}\ge\delta+\omega
 \ge\min\{\lambda_R,\delta+\omega\}.
\]
For condition (c), every $P_\xi(H_i)$ with $\xi<\delta+\omega$ is
infinite. Below $\delta$ it contains the infinite socle of the tail;
above $\delta$ use its finite-height socles. Conditions (d),(e) are
vacuous for the empty tuple. Therefore
\[
 H_i\le_{2\gamma+1}R.
\]
Lemma~\ref{lem:products} works at
every ordinal, so $H_i\oplus D\le_{2\gamma+1}R\oplus D$.
Every $\Sigma_{2\gamma+1}$ sentence true of $G$ is true in both
$H_i\oplus D$. Their reduced parts have different lengths, so at
least one is nonisomorphic to $G$. This proves the first assertion.

Suppose now that $\lambda_R<\omega\gamma$. If $\gamma=\beta+1$,
then $G_{\omega\beta}$ stabilizes after finitely many $p$-steps.
Such a terminal group has a $\Pi_3$ Scott sentence: finite-stage
stabilization, its finite-height Ulm invariants, and the rank of its
divisible part can all be specified at that level. Relativization to
$G_{\omega\beta}$, a $\Pi_{2\beta}$ domain, gives a
$\Pi_{2\beta+3}=\Pi_{2\gamma+1}$ sentence. The lower invariant
conditions have complexity at most $\Pi_{2\beta+1}$, so they fit.
The ambient Abelian $p$-group axioms are included.

If $\gamma$ is a limit, choose $\beta<\gamma$ and $n<\omega$ with
$\lambda_R=\omega\beta+n$. The same finite-stabilization argument
gives a $\Pi_{2\beta+3}$ Scott sentence; now
$2\beta+3<2\gamma$. Thus in both cases there is a
$\Pi_{2\gamma+1}$ Scott sentence. If a
$\Sigma_{2\gamma+1}$ Scott sentence also exists, Miller's theorem
in AGHT, Theorem~2.5, gives a $d\text{-}\Sigma_\eta$ Scott sentence
for some $\eta<2\gamma+1$ (in particular one at $2\gamma$).
Hence pure $\Sigma_{2\gamma+1}$ cannot be exact.
\end{proof}


\subsection{Deleting a direct summand of a bounded tail}
\begin{lemma}[Finite support absorption]
\label{lem:absorption}
Suppose $H=K\oplus T$ is bounded, $p^kT=0$, and a cyclic decomposition of
$K$ contains infinitely many $C_{p^k}$ summands. Then
$K\preceq_{\Sigma_1}H$, for ordinary infinitary formulas with parameters
from $K$. Consequently $H\le_2K$.
\end{lemma}

\begin{proof}
Existential truth passes from $K$ to $H$ by the subgroup inclusion.
Conversely, fix finite parameters in $K$ and finite witnesses in $H$.
Their $K$ coordinates and the parameters lie in a finite direct summand
$K_0$ of $K$. Their $T$ coordinates generate a finite group $T_0$.
Write $T_0$ as a finite sum of cyclic groups of orders $p^{\ell_i}$,
where $\ell_i\le k$. Choose distinct $C_{p^k}$ coordinates of $K$
outside the support of $K_0$, with generators $e_i$, and send the
$i$-th cyclic generator of $T_0$ to $p^{k-\ell_i}e_i$.
Together with the identity on $K_0$, this embeds $K_0\oplus T_0$ into
$K$ and fixes the parameters. It preserves equations and inequations,
so the witnesses can be moved into $K$. This applies to any true
existential disjunct of an infinitary $\Sigma_1$ formula.

Hence $(K,\bar b)\le_1(H,\bar b)$ for every finite tuple $\bar b$ in
$K$. To verify $H\le_2K$, answer each such tuple by itself in $H$.
The level-$1$ requirement is the displayed comparison, and the
level-$0$ requirement follows as well.
\end{proof}

The exponent bound is essential: an order-$p^{k+1}$ cyclic group cannot
be embedded into a fresh order-$p^k$ coordinate. Every finite higher
summand is retained in the applications below.

\begin{lemma}[Deleting a terminal direct summand]
\label{lem:terminal-deletion}
Put $\delta=\omega\gamma$, where $\gamma\ge0$, and suppose $G$ is reduced
with nonzero bounded $H=G_\delta$. Write $H=K\oplus T$, assume $K$ has
the same exponent as $H$, and put $Q=G/T$, with quotient map $q$.
For $A=\langle\bar a\rangle$ finite, suppose $A\cap H\subseteq K$.
Then
\[
 u_Q(\xi)=u_G(\xi)\quad(\xi<\delta),\qquad Q_\delta\cong K,\qquad
 \lambda_Q=\lambda_G,
\]
and
\begin{equation}\label{eq:odd-deletion}
             (G,\bar a)\le_{2\gamma+1}(Q,q\bar a).
\end{equation}
If also $H\le_2K$, then
\begin{equation}\label{eq:even-deletion}
                         G\le_{2\gamma+2}Q.
\end{equation}
\end{lemma}

\begin{proof}
Lemma~\ref{lem:quotient-filtration} applies to every $T\le G_\delta$:
$Q_\xi=q(G_\xi)$ through $\delta$, and all lower Ulm quotients are
unchanged. Thus $Q_\delta=H/T\cong K$. Equal exponents give equal
lengths and ensure $Q$ is reduced.

The pointed argument extends the finite retained-summand construction to arbitrary
$K$. The map $q$ is injective on $A+K$: if $a+x\in T$, where $a\in A$
and $x\in K$, then $a\in A\cap H\subseteq K$, so $a+x\in K\cap T=0$.
An element of $A$ outside $H$ retains its height by taking inverse
images of the quotient filtration through $\delta$. For an element
of $A\cap H$, height inside $H$ equals height in its direct summand
$K$; adding $\delta$ shows the full height is also unchanged.
Zero retains height infinity.

Apply ACM Theorem~4.4(1) to left $G$, right $Q$, and the inverse
partial map induced by $q$. Lower invariants, lengths, generated
subgroups, and heights meet clauses (a), (b), (d), and (e).
For (c), lift witnesses of any relative $\Sigma_1$ formula from
$Q_\delta=qK$ to $K$. Injectivity on $A+K$ preserves every positive
and negative atomic formula, including equations involving parameters
outside $H$. This is the required transfer from $(Q,q\bar a)$ to
$(G,\bar a)$, proving \eqref{eq:odd-deletion}.

For \eqref{eq:even-deletion}, apply ACM Theorem~4.4(2) with empty
tuples. Lower invariants agree. In the next block, the invariants of
$Q$ are bounded by those of $G$, since $K$ is a direct summand of $H$.
Lengths agree below $\delta+\omega$. In clause (c), use $H\le_2K$:
every $\Sigma_2$ sentence true in $K\cong Q_\delta$ holds in $H$.
For empty tuples this is precisely the required relative clause.
The subgroup and height clauses are vacuous.
\end{proof}

This is a sufficient algebraic criterion for the indicated pairs.
It is not an elimination of ACM's exceptional clauses for all pairs
and all distinguished tuples.

\section{Proof of the reduced classification}
\begin{proof}[Proof of Theorem~\ref{thm:reduced-classification}]
Put $\delta=\omega\gamma$. By \cite[Proposition~2.2]{ACM}, membership in $G_\delta$ is $\Pi_{2\gamma}$-definable. Specifying the
Ulm invariants below $\delta$ costs at most $\Pi_{2\gamma+1}$:
at $\omega\eta+r<\delta$, a dimension-at-least-$m$ condition is
$\Sigma_{2\eta+2}$, and an infinite dimension costs
$\Pi_{2\eta+3}$, where $\eta<\gamma$.

For finite nonzero $H$ and $\gamma>0$, require a
complete group table on $|H|$ distinct points of $G_\delta$ by a
$\Sigma_{2\gamma+1}$ formula, and bounds the size of $G_\delta$ by a
$\Pi_{2\gamma+1}$ formula. Along with lower invariants and the ambient
Abelian $p$-group axioms, this gives $d\text{-}\Sigma_{2\gamma+1}$.

With exactly one infinite multiplicity, write
\[
                       H=F_-\oplus C_{p^k}^{(\omega)}\oplus F_+,
\]
where $F_-$ and $F_+$ are finite sums of shorter and longer cyclic
groups. Theorem~\ref{thm:bounded-pgroup-classification} gives a $\Pi_2$ Scott sentence when $F_-=0$
and a $d\text{-}\Sigma_2$ Scott sentence otherwise. Relativization to
$G_\delta$ yields $\Pi_{2\gamma+2}$ and
$d\text{-}\Sigma_{2\gamma+2}$ respectively. A relative existential
block followed by a relative $\Pi_1$ matrix is
$\Sigma_{2\gamma+2}$; the dual universal form is
$\Pi_{2\gamma+2}$. The lower invariant conditions fit either bound.

With at least two infinite terminal multiplicities, ACM
Theorem~3.1(ii) gives $\Pi_{2\gamma+3}$. With length exactly $\delta$,
Theorem~3.1(iii) gives $\Pi_{2\gamma+1}$. Each description forces the
correct bounded tail, hence reducedness, and the complete Ulm
sequence. Ulm classification gives countable categoricity.
A nontrivial finite group ($\gamma=0$ and finite $H$) instead uses its
finite table and cardinality bound directly, giving
$d\text{-}\Sigma_1$ without adjoining higher-complexity ambient axioms.

\medskip\noindent\emph{Lower bounds for infinite terminal subgroups.}

The finite-$K$ instance of Lemma~\ref{lem:terminal-deletion} gives no
$\Sigma_{2\gamma+2}$ Scott sentence for any reduced $G$ with infinite
bounded $G_\delta$. It is the finite-$K$ case of
\eqref{eq:odd-deletion}: choose a finite direct summand containing
the parameters' intersection with $H$ and a maximal-order coordinate.
This also excludes every lower class below $\Pi_{2\gamma+2}$.
Thus the $\Pi_{2\gamma+2}$ row with $F_-=0$ is exact.

If $F_-\ne0$, take $T=F_-$ and
$K=C_{p^k}^{(\omega)}\oplus F_+$. The absorption lemma gives
$H\le_2K$, and \eqref{eq:even-deletion} gives
\[
                            G\le_{2\gamma+2}G/F_-.
\]
A nonzero shorter terminal invariant has been removed, so the groups
are nonisomorphic. Any $\Pi_{2\gamma+2}$ Scott sentence would transfer
to this nonisomorphic quotient. The upper bound and the pointed quotient
$\Sigma_{2\gamma+2}$ obstruction therefore give exact
$d\text{-}\Sigma_{2\gamma+2}$.

Suppose now that $m_i=m_k=\omega$ for $i<k$, and fix any finite
$\bar a$ in $G$. In a cyclic decomposition of $H$, retain the finitely
many $C_{p^i}$ coordinates supporting $A\cap H$, where
$A=\langle\bar a\rangle$. Let $T$ be the sum of all the other
$C_{p^i}$ coordinates, and let $K$ retain all remaining coordinates,
including every coordinate of order different from $p^i$. Then
\[
 H=K\oplus T,\qquad A\cap H\subseteq K,\qquad
 p^iT=0,\qquad C_{p^k}^{(\omega)}\text{ is a direct summand of }K.
\]
The exponent is unchanged. The two lemmas give, for $Q=G/T$,
\[
 (G,\bar a)\le_{2\gamma+1}(Q,q\bar a),\qquad
 G\le_{2\gamma+2}Q.
\]
The invariant $u_Q(\delta+i-1)$ is finite whereas
$u_G(\delta+i-1)=\omega$, so $Q\not\cong G$.

If a $d\text{-}\Sigma_{2\gamma+2}$ Scott sentence were $S\wedge P$,
choose $\bar a$ witnessing a true existential disjunct of its
$\Sigma_{2\gamma+2}$ part $S$, with $\Pi_{2\gamma+1}$ matrix.
The pointed comparison transfers this matrix to $(Q,q\bar a)$;
the unpointed comparison transfers its $\Pi_{2\gamma+2}$ part $P$.
This contradicts categoricity. Hence no such difference Scott
sentence exists. The $\Pi_{2\gamma+3}$ upper bound is exact.
An additional $\Sigma_{2\gamma+3}$ Scott sentence would, by Miller
collapse (AGHT Theorem~2.5 \cite{AGHT}), yield a lower difference
description and is also impossible.

\medskip\noindent\emph{Finite tails and exact limit lengths.}

Suppose $\gamma>0$ and $0<|G_\delta|<\omega$. Theorem~\ref{thm:odd-obstruction}
excludes a $\Sigma_{2\gamma+1}$ Scott sentence because the reduced
length is at least $\delta$. Put $Q=G/G_\delta$. The quotient lemma
preserves the invariants below $\delta$ and gives length $\delta$.
To verify the latter equality, if $Q_\xi=0$ for some $\xi<\delta$,
then $G_\xi\subseteq G_\delta$. If $p^mG_\delta=0$, this would give
$G_{\xi+m}=0$, where $\xi+m<\delta$, contradicting the length of $G$.
ACM Theorem~4.4(1), with left $G$, right $Q$, and empty parameters,
applies: the zero tail $Q_\delta$ embeds in $G_\delta$, so its
$\Sigma_1$ theory transfers in the required direction for clause (c).
Therefore
\[
                             G\le_{2\gamma+1}Q.
\]
Since $Q\not\cong G$, this excludes $\Pi_{2\gamma+1}$ and proves the
finite-terminal difference row. For nontrivial finite groups,
existential truth passes to a proper larger group and universal truth
passes to the trivial subgroup. Thus neither $\Sigma_1$ nor $\Pi_1$
is a Scott complexity, completing the $\gamma=0$ finite case.

Finally suppose $\lambda_G=\delta=\omega\gamma$, where $\gamma>0$.
Lemma~\ref{lem:prescribed-tail} supplies a reduced $B$ with the same
invariants below $\delta$ and $B_\delta\cong C_p$. Its length is
$\delta+1$, so it is nonisomorphic to $G$. We have
\[
               G\le_{2\gamma}B\qquad\hbox{and}\qquad B\le_{2\gamma}G.
\]
If $\gamma=\eta+1$, apply ACM Theorem~4.4(2) at index $\eta$ in
both directions. The invariants agree below
$\delta=\omega\eta+\omega$, and both lengths are at least $\delta$.
Every socle below $\delta$ in each group is infinite: a finite
socle in a reduced $p$-group makes the corresponding subgroup finite,
forcing the full length below the next limit $\delta$.
The finite-socle fact is stated just after ACM Theorem~4.4 and also
has a short direct proof. If a group with finite socle had unbounded
exponent, some nonzero socle element would have roots of arbitrarily
large finite depth. Each $p$-root fiber has at most the finite size of
the socle. K\"onig's lemma would therefore give an infinite compatible
root chain, generating a Pr\"ufer subgroup and contradicting
reducedness. Once the exponent is bounded by $p^m$, induction through
the $p$-map, whose fibers are finite, proves that the group is finite.
Thus the algebraic infinite-socle alternative in clause (c) applies
on both sides. Empty parameters meet the remaining clauses.
If $\gamma$ is a nonzero limit, $2\gamma=\gamma$ and ACM
Theorem~4.4(3) applies directly.

Consequently $G$ has no $d\text{-}\Sigma_{2\gamma}$ Scott sentence.
The $\Pi_{2\gamma+1}$ upper bound and Miller collapse prove exactness
for successor as well as limit $\gamma$. The trivial group is
described by $\forall x\,(x=0)$ and has no quantifier-free Scott
sentence. This proves every row of the classification.

\end{proof}
\section{Explicit successor realizations}\label{sec:walker}
For a countable ordinal $\xi$, let $W_\xi$ be the Abelian group with a
generator $e_s$ for each nonempty finite strictly decreasing sequence
\[
 s=(\xi,\xi_1,\ldots,\xi_n),\qquad
 \xi>\xi_1>\cdots>\xi_n,
\]
including the one-term sequence $(\xi)$. Its defining relations are
\[
 pe_{(\xi)}=0,
 \qquad pe_{s\mathbin{\smallfrown}\eta}=e_s
 \quad(\eta<\operatorname{last}(s)).
\]
This is the usual decreasing-sequence presentation of a Walker group;
the properties needed here follow directly from the presentation.

\begin{lemma}\label{lem:walker-filtration}
Every element of $W_\xi$ has a unique finite expression
\[
 \sum_s c_s e_s,\qquad 0\le c_s<p.
\]
Moreover, for every ordinal $\alpha$,
\begin{equation}\label{eq:walker-filtration}
 (W_\xi)_\alpha
 =\left\langle e_s:\operatorname{last}(s)\ge\alpha\right\rangle.
\end{equation}
Consequently,
\[
 (W_\xi)_\xi=\langle e_{(\xi)}\rangle\cong C_p,
 \qquad (W_\xi)_{\xi+1}=0.
\]
For every ordinal $\delta$,
\[
 (W_{\delta+1})_\delta\cong C_{p^2}.
\]
\end{lemma}

\begin{proof}
A finite ancestor-closed set of sequences carries the finite
presentation obtained by restricting the displayed relations to those
sequences. Adding a new leaf $s$ to such a finite set adjoins a root of
its parent with relation $pe_s=e_{s^-}$. The old finite group embeds in
this extension: in
\((H\oplus\mathbb Z)/\langle(-e_{s^-},p)\rangle\), the relation subgroup
meets $H\oplus\{0\}$ only at zero. The new quotient has exactly $p$
cosets over the old group. The root itself starts with a cyclic group
of order $p$. Induction over the finite tree therefore gives a unique
expression with all coefficients between $0$ and $p-1$.
All finite ancestor-closed presentations embed in larger ones by the
same argument, so the result holds in their directed union $W_\xi$.
In particular, there are no further relations between digit-normal
expressions. Each generator has finite $p$-power order, so the group is
indeed a countable Abelian $p$-group.

Let $S_\alpha$ be the set of nodes with last entry at least $\alpha$.
It is ancestor-closed. An element belongs to $\langle S_\alpha\rangle$
exactly when its digit-normal support is contained in $S_\alpha$.
At a successor, the $p$-multiples of generators in $S_\alpha$ lie in
$\langle S_{\alpha+1}\rangle$. Conversely, if
$\operatorname{last}(s)\ge\alpha+1$, the child
$s\mathbin{\smallfrown}\alpha$ exists, belongs to $S_\alpha$, and is
a $p$-root of $e_s$. Thus
\[
 p\langle S_\alpha\rangle=\langle S_{\alpha+1}\rangle.
\]
At a limit $\lambda$, uniqueness of finite digit support gives
\[
 \bigcap_{\alpha<\lambda}\langle S_\alpha\rangle
 =\langle S_\lambda\rangle.
\]
Transfinite induction proves \eqref{eq:walker-filtration}.
At level $\xi$ only the root remains. At level $\delta$ in
$W_{\delta+1}$, exactly the root and its child $(\delta+1,\delta)$
remain; the latter has $p$-multiple equal to the former, giving a
cyclic group of order $p^2$.
\end{proof}

Transfinite $p$-subgroups commute with direct sums:
\begin{equation}\label{eq:p-filtration-direct-sums}
 \left(\bigoplus_i G_i\right)_\alpha
 =\bigoplus_i(G_i)_\alpha.
\end{equation}
The successor step is immediate. At a limit, an element has finite
coordinate support, and it belongs to each previous subgroup exactly
when each coordinate does. This proves the identity by induction.
The corresponding Ulm invariants add.

For $\delta=\omega\gamma$, define
\begin{align}
 B_\delta&=\bigoplus_{\xi<\delta}W_\xi^{(\omega)},
 \label{eq:lower-reservoir}\\
 A^{\rm odd}_{\gamma,r}
 &=B_\delta\oplus W_\delta^r,
 \label{eq:odd-group}\\
 A^{\rm even}_{\gamma,r}
 &=B_\delta\oplus W_\delta^r
                 \oplus W_{\delta+1}^{(\omega)}.
 \label{eq:even-group}
\end{align}
These are countable presentations for every countable $\gamma$.
They need not be computable presentations when $\gamma$ is
noncomputable; no effectivity of arbitrary countable ordinal input is
asserted.

Each copy of $W_\xi$ contributes one dimension to its $\xi$-th Ulm
invariant. The infinitely many copies in \eqref{eq:lower-reservoir}
therefore give
\[
 u_{B_\delta}(\xi)=\omega\quad(\xi<\delta),
 \qquad (B_\delta)_\delta=0.
\]
Lemma~\ref{lem:walker-filtration} now gives the complete invariant
lists needed for the two families:
\begin{center}
\begin{tabular}{c|c|c|c|c}
Group & $\xi<\delta$ & $\xi=\delta$ & $\xi=\delta+1$ & later $\xi$\\
\hline
$A^{\rm odd}_{\gamma,r}$ & $\omega$ & $r$ & $0$ & $0$\\
$A^{\rm even}_{\gamma,r}$ & $\omega$ & $r$ & $\omega$ & $0$
\end{tabular}
\end{center}
Their lengths are $\delta+1$ and $\delta+2$, respectively, and
\begin{equation}\label{eq:terminal-groups}
 (A^{\rm odd}_{\gamma,r})_\delta\cong C_p^r,
 \qquad
 (A^{\rm even}_{\gamma,r})_\delta
 \cong C_p^r\oplus C_{p^2}^{(\omega)}.
\end{equation}
Both are reduced because their transfinite $p$-filtrations reach zero.
For different $r$, their $\delta$-th Ulm invariants differ.


\begin{theorem}[The difference spectrum]\label{thm:difference-spectrum}
Every positive countable successor $d\text{-}\Sigma_\alpha$ complexity
occurs among countable reduced abelian $p$-groups. Among infinite such
groups the attained difference levels are exactly the successors
$\alpha\geq2$.
\end{theorem}
\begin{proof}
For $\gamma>0$ the two groups in \eqref{eq:odd-group} and
\eqref{eq:even-group} have terminal groups displayed in
\eqref{eq:terminal-groups}. Theorem~\ref{thm:reduced-classification} gives
exact complexities $d\text{-}\Sigma_{2\gamma+1}$ and
$d\text{-}\Sigma_{2\gamma+2}$ respectively. At level two take
$C_p^r\oplus C_{p^2}^{(\omega)}$, $r\geq1$; nontrivial finite groups
give level one. Every successor ordinal is $2\gamma+1$ or $2\gamma+2$:
write it as a limit (or zero) plus a positive finite remainder and use
$2\lambda=\lambda$ for nonzero limit $\lambda$.
The general nonzero-limit exclusion is the published theorem
\cite[Theorems~1.6 and 2.6]{AGHT}, not a new restriction specific to
$p$-groups.
Finally, an infinite $p$-group cannot have a $d\text{-}\Sigma_1$ Scott
sentence. A true existential disjunct of its $\Sigma_1$ part has finite
witnesses. They generate a finite subgroup, to which the universal
part passes. This would be a nonisomorphic countable model of the same
sentence.
\end{proof}

In the rank conventions of \cite{GRT}, exact
$d\text{-}\Sigma_{\beta+1}$ complexity gives
$\operatorname{SR}_p=\beta<\operatorname{SR}=\beta+1$.
The displayed families therefore give explicit rank-gap examples at
every countable $\beta\geq1$. This supplies concrete examples for their
Question~1; their prior abstract existence result remains credited.
No uniform effective presentation from an arbitrary countable ordinal
notation is being asserted.

\section{Arbitrary divisible parts}
\label{sec:infinite-divisible}

The final even-form case can be treated without assuming that ACM's
reduced-group statement automatically applies to a nonreduced group.
We prove the sufficient fragment that is needed. Give each nonzero
element of the divisible part height $\infty$, larger than every ordinal;
as usual $h(0)=\infty$ as well. A group with infinitely many independent
Pr\"ufer summands has $P_\xi$ infinite for every ordinal $\xi$.

\begin{lemma}[Infinite-divisible-rank comparison]
\label{lem:infinite-divisible-bf}
Let $A,B$ be countable Abelian $p$-groups whose divisible ranks are
both infinite. Let $g:\langle\bar b\rangle\to\langle\bar a\rangle$
be the isomorphism induced by the indicated finite tuples.
The following conditions are sufficient for the displayed comparisons.
\begin{enumerate}
\item For $\delta=\omega\gamma$,
$(A,\bar a)\le_{2\gamma+1}(B,\bar b)$ provided the Ulm
invariants agree below $\delta$ and, for $a=g(b)$, either
\[
 h_A(a)=h_B(b)<\delta,
 \quad\text{or}\quad
 h_B(b)\ge\delta,\qquad
 h_A(a)\ge\min\{h_B(b),\delta+\omega\}.
\]
\item $(A,\bar a)\le_{2\gamma+2}(B,\bar b)$ provided the
Ulm invariants agree below $\delta=\omega\gamma$,
\[
 u_B(\xi)\le u_A(\xi)
 \quad(\delta\le\xi<\delta+\omega),
\]
and each $a=g(b)$ either has the same height as $b$ below
$\delta+\omega$, or both heights are at least $\delta+\omega$.
\item At a nonzero limit $\eta$, $(A,\bar a)\le_\eta(B,\bar b)$
provided the Ulm invariants agree below $\omega\eta$ and heights
are equal below $\omega\eta$ and simultaneously above that cut.
\end{enumerate}
\end{lemma}

\begin{proof}
We verify the extension argument behind the infinite-socle cases of
ACM, Theorem~4.4, rather than importing its reducedness hypothesis.
The following algebraic ingredients apply to arbitrary Abelian
$p$-groups. For a finite subgroup $S$, write $S_\xi=S\cap G_\xi$.

If $S$ is finite and $x\notin S$ has $px\in S$, one may choose
its representative modulo $S$ to be proper: its height is maximal
in its finite coset. Among such representatives also maximize the
height of its $p$-multiple. These two maxima exist because the coset
is finite, even when the maximum is $\infty$. Then
$h(x+s)=\min\{h(x),h(s)\}$ for every $s\in S$.
If its height is $\xi$ and $h(px)>\xi+1$, a root $z$ of $px$
can be chosen in $G_{\xi+1}$, and $x-z$ is a proper socle
element at height $\xi$. The available dimension is measured by
the elementary map
\[
 S_\xi^*/S_{\xi+1}\longrightarrow P_\xi/P_{\xi+1},
 \qquad [s]\longmapsto[s-z],
\]
where $S_\xi^*=S\cap G_\xi\cap p^{-1}G_{\xi+2}$ and
$z\in G_{\xi+1}$ satisfies $pz=ps$. It is injective and
well-defined. A proper socle element can be chosen precisely when
its image is not the entire Ulm quotient. Thus equality, or the
appropriately directed inequality, of Ulm dimensions permits the
usual extension at a specified ordinal height. This is the
argument of ACM, Lemma~2.3; it requires no reducedness.

Finally, at a height where only a lower bound is required, infinitely
many socle elements permit avoidance of the finitely many forbidden
choices caused by $S$. Here this is available at every height from
the infinite divisible part. Roots of a divisible element can be
chosen divisible. More generally, to obtain a root of height at
least $\rho$, the only assertion needed is that the given element
lies in $G_{\rho+1}$. No predecessor of an infinite or limit height
is taken.

Induct on the comparison index. At index zero, the given finite
group isomorphism is the required atomic agreement. For an odd
index $2\gamma+1$, first fix the entire finite right extension $B'$.
Choose a positive integer $M$ above all $k$ for which an element
of $B'$ has height $\delta+k$. Let $r$ be the number of cyclic
one-element extensions needed and put $N=M+r+1$. Build $B'$ over
$\langle\bar b\rangle$ one element at a time, using the proper
representatives with the secondary maximum just specified.
After $i$ extensions, maintain equality of heights below $\delta$;
for a right element of finite block height $\delta+k$ require its
image to have height at least $\delta+k$; and for a right element
of height at least $\delta+\omega$ require its image to have height
at least $\delta+N-i$. The original odd condition implies this
finite-headroom invariant at $i=0$.

Below $\delta$, a one-step height has a root of exactly the same
height, and the secondary maximum prevents that root from lying
in the preceding finite image. A height jump is handled by the
Ulm-quotient dimension calculation above. In detail, in the one-step
case put $\xi=h_B(x)$ and $y=g(px)$. A root $w$ of $y$ has
height exactly $\xi$. If $w+g(s)$ had greater height for some
$s$ in the preceding subgroup, then $x+s$ would still be a proper
representative of height $\xi$. Moreover,
$p(w+g(s))=g(p(x+s))$ would belong to $A_{\xi+2}$.
Since $\xi+2<\delta$ and the old map preserves membership at
that cut, $p(x+s)\in B_{\xi+2}$. Its height would therefore
exceed $h_B(px)=\xi+1$, contrary to the secondary maximum.
Thus $w$ is proper over the preceding image. This also rules out
$w$ in the preceding image. In the jump case choose a root $w_1$
of $y$ in $A_{\xi+1}$ and a proper $w_2\in P_\xi(A)$, using
the unused Ulm-quotient dimension; $w=w_1+w_2$ is proper of height
$\xi$. In either case the map extends on all sums by the
proper-coset identity.

For a finite block height
$\delta+k$, the image of the $p$-multiple has height at least
$\delta+k+1$, so a root of sufficient height exists; infinitely
many socle elements of at least that height permit avoidance of the finite
image. For a right height at least $\delta+\omega$, use a root
of height at least $\delta+N-i-1$, again adding an unused socle
element when necessary. The proper-coset identity proves the
invariant for sums whose preceding summand has a different
right height; when both heights are above $\delta$, only the
indicated lower bound is needed. Thus the invariant holds on the
entire new finite subgroup. At completion, all heights are equal
when truncated at $\delta$.

In particular, this step does not demand a root still above the
limit $\delta+\omega$ of an element of height exactly that limit:
only the explicitly finite headroom is spent.

If $\gamma=\beta+1$, the reversed map now satisfies part (2) at
index $2\beta+2=2\gamma$: all invariants below $\delta$ agree and
heights are equal below that cut. If $\gamma$ is a limit, it
satisfies part (3) at $2\gamma=\gamma$. For $\gamma=0$, it
satisfies the atomic condition. The induction hypothesis therefore
gives the response required by the definition of $\le_{2\gamma+1}$.

At an even index $2\gamma+2$, first fix the entire finite right
extension $B'$ and put
\[
 M=1+\max\bigl(\{k<\omega:
       h_B(x)=\delta+k\text{ for some }x\in B'\}\cup\{0\}\bigr).
\]
Let $r$ be the number of one-element extensions needed and put
$N=M+r+1$. During the $i$-th extension maintain equal heights
below $\delta+M$, and send every element of right height at least
$\delta+\omega$ to left height at least $\delta+N-i$.
The initial map satisfies this invariant. Below $\delta+M$,
root and jump extensions use equality of the lower invariants and
$u_B\le u_A$ in the final block. The root/properness argument is
the same as above. The Ulm-quotient calculation also applies at
the greatest possible finite height $\delta+M-1$: source elements
above that height are already at least $\delta+\omega$, because
there are no right-hand heights in
$[\delta+M,\delta+\omega)$ anywhere in $B'$.
Before every extension $N-i\ge M+2$, and after all $r$ extensions
$N-r=M+1$. Thus the old map preserves membership through
$\xi+2\le\delta+M+1$, including at the boundary
$\xi=\delta+M-1$. This is the exact membership preservation
used in the preceding root/properness argument as well.
Hence the partial map preserves the predicates needed
for $S_\xi^*$ and $S_{\xi+1}$ at every height used in the calculation.
Their dimensions agree in the two finite subgroups, and the stated
Ulm inequality supplies the unused target class whenever one is
needed in the source.

For a right element of height
at least $\delta+\omega$, its $p$-multiple's image has a root
of height at least $\delta+N-i-1$; an unused socle element of at
least that height avoids the preceding finite subgroup. Every right
element of $B'$ has either a height below $\delta+M$ or a height
at least $\delta+\omega$, by the definition of $M$.
For a preceding summand of height below $\delta+M$, its height
and that of its image agree and are below the new root's height;
the sum has that same smaller height. For a preceding summand
above the cut, both image heights meet the new headroom bound,
so their sum does too. This proves the invariant on all sums;
nonzero coefficients modulo $p$ do not change heights. At the end,
the inverse map satisfies the odd
height condition at $2\gamma+1$. Its lower invariants agree;
its left-hand socles are all infinite. The odd induction step
supplies the required comparison of the extended tuples.

At a limit $\eta$, each fixed smaller successor index uses a cut
strictly below $\omega\eta$. The assumed equality of invariants
and truncated heights therefore supplies its sufficient conditions.
Apply the induction hypothesis at a successor index above the
requested smaller index to obtain the required finite response.
This proves part (3) and completes the induction.

The only uses of the infinite-socle assumption in these extensions
are the explicitly available avoidance choices. In particular, the
reduced-case inference from length to infinite socles is unnecessary,
and no exceptional finite-socle alternative is invoked.
\end{proof}

\begin{theorem}[No exact pure-$\Sigma$ complexity for any countable Abelian $p$-group]
\label{thm:no-pure-sigma}
No countable Abelian $p$-group has exact pure
$\Sigma_\alpha$ Scott complexity at any countable index.
\end{theorem}

\begin{proof}
At index zero, every closed group-language term equals $0$, so
all Abelian $p$-groups have the same quantifier-free sentences.
Consequently none has a $\Sigma_0$ Scott sentence.
Write $G=R\oplus D$. General nonzero-limit and low-index exclusions, and
Theorem~\ref{thm:odd-obstruction}, leave only even-form
indices $2\gamma+2$. Put $\delta=\omega\gamma$.

If $\lambda_R\ge\delta+\omega$, the odd obstruction at
$\gamma+1$ excludes a $\Sigma_{2\gamma+3}$ Scott sentence and
therefore a $\Sigma_{2\gamma+2}$ one. If $\lambda_R<\delta$,
the short-length upper bound is already $\Pi_{2\gamma+1}$.
If $R_\delta$ is infinite of bounded exponent, choose for each
finite reduced-coordinate tuple a finite retained summand as in
Lemma~\ref{lem:terminal-deletion}. Its pointed odd comparison,
combined with Lemma~\ref{lem:products}, transfers the matrix of any
putative $\Sigma_{2\gamma+2}$ Scott sentence to a nonisomorphic quotient
with the same divisible part.

It remains that $\delta\le\lambda_R<\delta+\omega$ and
$F=R_\delta$ is finite, possibly zero. Suppose first that $D$ has
finite rank. The group $F\oplus D$ has a $\Pi_2$ Scott sentence.
To see this, choose $n$ killing $F$ and impose stabilization
$p^nH=p^{n+1}H$, a $\Pi_2$ sentence. For each $k\le n$, fix the
finite dimension of $(p^kH)[p]$. Having dimension at least $r$ is
$\Sigma_1$, witnessed by independent elements $p^kx_i$ with
$p^{k+1}x_i=0$; exact finite dimension is $d\text{-}\Sigma_1$.
These dimensions specify all finite Ulm invariants and the finite
divisible rank. Along with the Abelian $p$-group axioms they are
countably categorical. Relativizing this $\Pi_2$ sentence to the
$\Pi_{2\gamma}$ domain $G_\delta$ gives a
$\Pi_{2\gamma+2}$ Scott sentence after the lower Ulm list is
included. If a $\Sigma_{2\gamma+2}$ Scott sentence exists too,
Miller's theorem lowers it to $d\text{-}\Sigma_{2\gamma+1}$.
Thus pure $\Sigma_{2\gamma+2}$ is not exact.

Suppose instead that $D$ has infinite rank. Replace the finite
terminal Ulm factor $F$ by $F\oplus C_p$, leaving every factor
before $\gamma$ unchanged. If $F=0$, append the new final factor
$C_p$. The earlier factors are unbounded because
$\lambda_R\ge\delta$. For $\gamma>0$, Zippin's theorem supplies a reduced
group $R'$ with this sequence; for $\gamma=0$ take $R'=F\oplus C_p$ directly. Put $G'=R'\oplus D$.
The Ulm invariants of $G'$ and $G$ agree below $\delta$; in the
block from $\delta$ to $\delta+\omega$, the former dominate the
latter, with strict inequality at $\delta$. Both groups have
infinite divisible rank. Lemma~\ref{lem:infinite-divisible-bf}(2),
with left group $G'$ and right group $G$ and empty tuples, gives
\[
 G'\le_{2\gamma+2}G.
\]
Every $\Sigma_{2\gamma+2}$ sentence true of $G$ is therefore true
of $G'$, whereas their Ulm invariants differ. Hence $G$ has no
$\Sigma_{2\gamma+2}$ Scott sentence in this last case.
All cases are exhausted.
\end{proof}



\section{Answers to the source questions and remaining work}
Theorem~\ref{thm:reduced-classification} gives a negative answer to
\cite[Question~5.1]{ACM}. Infinite Ulm invariants below the terminal
block do not force the generic upper bound to be optimal. For a concrete
example use
\[
 G=B_\omega\oplus W_\omega^{(\omega)}.
\]
Its finite-height Ulm invariants and its $\omega$th invariant are all
$\omega$, and its length is $\omega+1$. Nevertheless its terminal group
is $C_p^{(\omega)}$, so its exact complexity is $\Pi_4$, whereas the
generic case-(ii) bound of \cite[Theorem~3.1]{ACM} is $\Pi_5$.
More precisely, that generic bound is optimal exactly when at least two
terminal cyclic multiplicities are infinite. At exact limit length
the $\Pi_{2\gamma+1}$ bound is always optimal; for nonzero finite tails,
the generic difference bound improves by one index.

Theorem~\ref{thm:no-pure-sigma} and
Theorem~\ref{thm:difference-spectrum} answer both realization alternatives
of \cite[Question~5.2]{ACM}: the pure-$\Sigma$ spectrum is empty, and the
difference spectrum consists of all positive countable successor levels.
The impossibility of exact limit difference levels is the general
published restriction, while the successor realizations already lie in
the reduced subclass.

The optimal complexity of each individual nonreduced group is not
classified here. Nor do the special comparison lemmas eliminate all
relative-sentence conditions from the general pointed back-and-forth
analysis requested in \cite[Question~5.3]{ACM}. First-order Vaught's
conjecture is not resolved by these Scott-sentence results.

\paragraph{Preparation and status.}
This manuscript was prepared with AI assistance from a research programme
directed by the repository owner. The assembly combines the argument
for successor realizations, its transfinite-existence repair, the
divisible-rank argument, and the terminal classification into one
canonical text. Published dependencies are credited at use. The
unbounded assertions rest on the displayed hand proofs and those
published inputs; no independent human audit, formal verification, or
historical priority claim is asserted. No private notes, named project
claims, or conventions outside this repository are required.

\begin{thebibliography}{9}
\bibitem{ACM} R.~Alvir, B.~F.~Csima and L.~MacLean,
\emph{Scott Complexity of Reduced Abelian $p$-Groups},
arXiv:2407.06940v1 (2024).
\url{https://arxiv.org/abs/2407.06940v1}.
\bibitem{AGHT} R.~Alvir, N.~Greenberg, M.~Harrison-Trainor and
D.~Turetsky, \emph{Scott complexity of countable structures},
Journal of Symbolic Logic 86(4) (2021), 1706--1720.
\url{https://doi.org/10.1017/jsl.2021.4}.
\bibitem{CT} F.~Calderoni and S.~Thomas,
\emph{The bi-embeddability relation for countable abelian groups},
Transactions of the American Mathematical Society 371(3) (2019),
2237--2254. Zippin existence is used in the form of Theorem~4.2,
printed p.~10 of the author manuscript.
\url{https://sites.math.rutgers.edu/~sthomas/TFA.pdf}.
\bibitem{GRT} D.~Gonzalez, D.~Rossegger and D.~Turetsky,
\emph{Scott Analysis below the Vaught Ordinal},
arXiv:2606.15205v2 (2026), Question~1 and the Scott-rank conventions.
\url{https://arxiv.org/abs/2606.15205v2}.
\end{thebibliography}
\end{document}
